\documentclass[12pt, a4paper]{report}
\usepackage{amsmath}
\usepackage[english]{babel}
\usepackage{graphicx}
\usepackage{times}
\usepackage{geometry}
\usepackage{csquotes}
\usepackage{parskip}
\usepackage{fancyhdr}
\usepackage{lipsum}
\usepackage{graphicx}
\usepackage{hyperref}
\usepackage{longtable}
\usepackage{pgfplots}
\pgfplotsset{compat=1.18}
\usepackage[style = apa, backend = biber]{biblatex}
\usepackage{setspace}
\usepackage{array}
\usepackage{float}
\usepackage{booktabs}
\usepackage{multicol}
\usepackage{esdiff}
\usepackage{titlesec}  % Importa il pacchetto titlesec per personalizzare lo stile del capitolo
\usepackage{titletoc}
\usepackage{subcaption}

% Imposta il formato dei capitoli sulla stessa riga senza salto pagina
\titleformat{\chapter}[hang] 
  {\normalfont\LARGE\bfseries} % Formato testo
  {Chapter \thechapter:} % Mostra "Chapter X:" 
  {1em} % Spazio tra "Chapter X:" e il titolo
  {} % Nessuna modifica aggiuntiva

% Rimuovi il salto di pagina con titlespacing
\titlespacing*{\chapter}{0pt}{-20pt}{20pt}  % Nessun salto pagina

% Configurazioni di default
\renewcommand{\thechapter}{\Roman{chapter}}
\renewcommand{\thesection}{\arabic{chapter}.\arabic{section}}
\renewcommand{\thesubsection}{\arabic{chapter}.\arabic{section}.\arabic{subsection}}
\renewcommand{\thetable}{\arabic{chapter}.\arabic{table}}
\renewcommand{\thefigure}{\arabic{chapter}.\arabic{figure}}
\renewcommand{\theequation}{\arabic{chapter}.\arabic{equation}}
\usepackage{caption}
\captionsetup[figure]{labelformat=simple, labelsep=none, labelfont={bf}, textfont={}, singlelinecheck=off}
\DeclareCaptionFormat{custom}{#1#2\\#3} 
\captionsetup[figure]{format=custom}
\captionsetup[table] {labelformat=simple, labelsep=none, labelfont={bf}, textfont={it}, singlelinecheck=off}
\renewcommand{\tablename}{Tabla}
\captionsetup[table]{format=custom}

% Recursos bibliográficos
\addbibresource{ref.bib}

% Formato di pagina
\geometry{a4paper, left=3.5cm,right=2.54cm,top=2.54cm,bottom=2.54cm}
\setlength{\parindent}{2em}
\pagestyle{fancy}
\fancyhf{} % clear all header and footer fields
\renewcommand{\headrulewidth}{0pt} % no line in header area
\fancyfoot[R]{\thepage} % page number in the right footer
\fancypagestyle{plain}{%
  \fancyhf{} % clear all header and footer fields
  \fancyfoot[R]{\thepage} % page number in the right footer
}
\begin{document}

% Frontespizio con Table of Contents
\begin{titlepage}
    \centering
    {\scshape\large University of Genova \par}
    \vspace{0.3cm}
    {\scshape\large Department of Computer Science, Bioengineering, Robotics and System Engineering (DIBRIS) \par}
    \vspace{0.3cm}
    \includegraphics[width=0.2\textwidth]{mainmatter/Images/Unige-logo.jpeg}\\[0.3cm]
    
    {\LARGE\bfseries COMPUTER VISION \par}
    \vfill
    {\scshape\large Lab Session n.5 \par}
    \vfill
    
    \begin{minipage}[t]{0.45\textwidth}
        \begin{flushleft} \large
            7852527 - Mattia Tinfena\\
            6504193 - Elisa Martinenghi\\
            7687956 - Chiara Buono
        \end{flushleft}
    \end{minipage}
    
    \vfill

\end{titlepage}

% Inizio del documento principale
\pagenumbering{arabic}  % Riprendi la numerazione delle pagine da qui
\chapter*{Introduction}


\chapter{NCC-based segmentation}



\chapter{Harris corner detection}


\newpage
\section*{Conclusions}





\end{document}
